\chapter{Layer L1a: Hepp trees}\label{ch:hepp}

The multiscale bookkeeping structure of the paper's \S5.2, on the plane
carrier with order-forgetting automorphisms
(\texttt{DESIGN.md}~\S5.5).

\begin{definition}[Hepp trees, markings, multiplicities; paper
  Defs.~5.1, 5.3]\label{D-hepp}
  \lean{Anderson4D.PlaneTree, Anderson4D.PlaneTree.HeppMarking,
    Anderson4D.PlaneTree.Multiplicities, Anderson4D.PlaneTree.Aut,
    Anderson4D.PlaneTree.card_aut_eq_autCard}\leanok
  A \emph{Hepp tree} is a plane rooted tree in which every internal node
  has at least two children (carrier \texttt{PlaneTree}; leaves are
  identified by their position paths), together with, as separate data:
  a \emph{marking} \(N\colon\)~branching nodes~\(\to 2^{\mathbb{N}}\),
  strictly increasing toward the root, and \emph{multiplicities}
  \(m\colon\)~leaves~\(\to \mathbb{N}_{\ge 2}\). Derived notation:
  \(\mathcal{B}\), \(\mathcal{L}\), \(\gamma_{\mathfrak n}\), subtrees,
  least common ancestors. Automorphisms are \emph{order-forgetting}:
  recursive child permutations preserving subtree shape (and marks, for
  \(\mathrm{Aut}(\mathcal{T},N)\)); concretely accessed through recursive
  canonicalization. The finite group \(\mathrm{Aut}(\mathcal{T})\), its
  action on leaves and markings, and the stabilizer description of
  \(\mathrm{Aut}(\mathcal{T},N)\) are part of the checked definition layer.
\end{definition}

\begin{definition}[Admissible embeddings; paper Def.~5.4]\label{D-adm}
  \uses{D-hepp}
  \lean{Anderson4D.IsAdmissible, Anderson4D.Realizes,
    Anderson4D.RealizesTuple}\leanok
  For a marked Hepp tree: injections \((z_{\mathfrak l})\colon
  \mathcal{L} \to \mathbb{Z}^4_M\) with
  \(\lvert z_{\mathfrak l} - z_{{\mathfrak l}'}\rvert \ge
  \tfrac12 N_{{\mathfrak l}\vee{\mathfrak l}'}\) and the linked-children
  condition (b); realizability \(Z \leftrightarrow (\mathcal{T}, N)\);
  the multiset/multiplicity variants.
\end{definition}

\begin{lemma}[Existence of realizing trees; paper Lemma~5.5]\label{L-5.5}
  \uses{D-adm}
  \lean{Anderson4D.rdec_exists_realizing_marked_tree,
    Anderson4D.exists_realizing_tree,
    Anderson4D.exists_realizing_tree_of_across_pairing}\leanok
  Every finite \(Z \subseteq \mathbb{Z}^4_M\) is realized by some marked
  Hepp tree with \(1 < N_{\mathfrak n} \lesssim M\) (greedy scale
  construction); every paired vector \((y_j)\) is realized in the
  multiplicity sense with even multiplicities.
\end{lemma}

\begin{lemma}[Sum decomposition; paper (5.6)--(5.11)]\label{R-decomp}
  \uses{L-5.5,L-5.2,D-pair,D-int,D-adm,D-ledger}
  \lean{Anderson4D.latticeChainSum_le_treeSum,
    Anderson4D.sum_eq_sum_paired_tree_incidence_div,
    Anderson4D.PositiveAssignment.card_le_two_pow,
    Anderson4D.pairedWordSum_le_paperSum,
    Anderson4D.primitivePairedWordSum_le_paperSumFiltered}\leanok
  The master reduction of the counting sum (5.5): fix the tree
  (\(C^n\) choices by \ref{L-5.2}), the multiplicities and the ordering
  data; the symmetry-factor denominator; and the word-sum bound (5.10)
  in the normalization of \ref{D-ledger}, with
  \(\mathrm{pairedWordSum}\) on the left and
  \(\big(\prod_{\mathfrak l}(m_{\mathfrak l}/2)!/m_{\mathfrak l}!\big)
  \mathrm{paperSum}\) on the right exactly as printed.
\end{lemma}

\begin{lemma}[Parent-function tree bound; used at paper (5.24)]\label{I-cayley}
  \lean{Anderson4D.parentCode_one_add_bound_real,
    Anderson4D.weightedCayley_le_of_parent_injection,
    Anderson4D.sum_fullParentWeight,
    Anderson4D.step5_parent_volume_bound}\leanok
  The sufficient parent-function injection estimate
  \(\sum_{\mathcal H}\prod_i (1 + w_i)^{d(i)-1}
  \le \big(q + \sum_i w_i\big)^{q-1}\) used in Step~5 of the proof of
  \ref{P-5.6}.  This replaces the paper's exact weighted Cayley identity
  by the weaker conclusion actually needed in (5.24)--(5.25), obtained
  from an injection into parent functions.
\end{lemma}

\begin{proposition}[Volume estimate; paper Prop.~5.6]\label{P-5.6}
  \uses{D-adm,L-5.2,I-cayley}
  \lean{Anderson4D.volume_estimate}\leanok
  The orbit-stabilizer lower bound (5.12) for the number of markings
  realizing a given embedding, and the embedding-counting upper bounds
  (5.13)--(5.14) with automorphism factors
  \(\lvert\mathrm{Aut}(\mathcal{T})\rvert /
  \lvert\mathrm{Aut}(\mathcal{T},N)\rvert\) and
  \(\prod_{\mathfrak n} N_{\mathfrak n}^{4(\gamma_{\mathfrak n}-1)}\)
  (and the \(\langle z_0 - z_1\rangle^{-4}\)-refined variant).
  Six-step proof as in the paper's \S5.3.
\end{proposition}
